\section{An explicit arbitrary-depth obstruction to finite Laguerre testing}
\label{sec:finitedepth}
The preceding Gaussian--hyperbolic example
shows that radial curvature does not force even
the first Fourier-side Laguerre inequality.
There is a stronger structural obstruction:
\emph{no fixed finite number of generalized
Laguerre inequalities, even when valid for
every real Fourier variable, implies the
real-zero property.}

\begin{theorem}[Finite-depth obstruction]
\label{thm:finitedepth}
For every integer $N\ge1$ put
\[
d_N=2^{-(N+2)},\qquad
K_N(u)=e^{-u^2/2}
\bigl(1+d_N e^{-1/2}\cosh u\bigr)
\]
and $F_N(z)=\int_{\mathbb R}K_N(u)e^{izu}\,du$.
Then $K_N$ is positive, even, rapidly
decreasing and
\begin{equation}
F_N(z)=\sqrt{2\pi}e^{-z^2/2}
\bigl(1+d_N\cos z\bigr).
\label{eq:finiteF}
\end{equation}
For every $1\le n\le N$ and $x\in\mathbb R$,
\begin{equation}
\boxed{
L_n[F_N](x)\ge
\frac{\pi e^{-x^2}}{n!}>0.
}
\label{eq:finitelb}
\end{equation}
Nevertheless $F_N$ has nonreal zeros
\begin{equation}
z=(2k+1)\pi\pm
i\operatorname{arcosh}(1/d_N),
\qquad k\in\mathbb Z.
\label{eq:finitezero}
\end{equation}
\end{theorem}

\begin{proof}
Gaussian Fourier integration gives
\eqref{eq:finiteF} and
\eqref{eq:finitezero} follows as in
Theorem~\ref{thm:coshcounter}.
The generalized Laguerre coefficients
are determined by the generating function
\begin{align*}
&F_N(x+iy)F_N(x-iy)\\
&\quad=2\pi e^{-x^2}e^{y^2}
\bigl[1+2d_N\cos x\cosh y\\
&\qquad\quad+\tfrac12d_N^2(\cosh 2y+\cos 2x)\bigr].
\end{align*}
Writing $L_n[F_N](x)$ for the
coefficient of $y^{2n}$, set
\begin{align*}
a_n&=\frac1{n!},\\
b_n&=\sum_{j=0}^n
  \frac1{(n-j)!(2j)!},\\
c_n&=\sum_{j=0}^n
  \frac{4^j}{(n-j)!(2j)!}.
\end{align*}
Then
\begin{align}\label{eq:finitecoefs}
\frac{L_n[F_N](x)}{2\pi e^{-x^2}}
&=a_n+2d_Nb_n\cos x\nonumber\\
&\quad+\frac{d_N^2}{2}
\left(c_n+a_n\cos 2x\right).
\end{align}
The last parenthesis is nonnegative
because $c_n\ge a_n$. Next
\[
\frac{b_n}{a_n}
=\sum_{j=0}^n
\binom nj\frac{j!}{(2j)!}
\le\sum_{j=0}^n\binom nj=2^n.
\]
Hence for $1\le n\le N$,
\begin{align*}
\frac{L_n[F_N](x)}{2\pi e^{-x^2}}
&\ge a_n-2d_Nb_n\\
&\ge a_n(1-2d_N2^n)\\
&\ge a_n/2>0,
\end{align*}
which proves \eqref{eq:finitelb}.
All bounds hold uniformly for
$x\in\mathbb R$, not just on compact sets.
\end{proof}

\begin{proposition}[First-order positivity hides
a second-order failure]
\label{prop:firstsecond}
With $d=2/5$, set
\[
F(z)=\sqrt{2\pi}\,e^{-z^2/2}
\left(1+\frac25\cos z\right).
\]
Then $F(x)>0$ and $L_1[F](x)>0$
for every real $x$, but
\begin{equation}
\boxed{L_2[F](\pi)=
-\frac{2\pi e^{-\pi^2}}{25}<0.}
\end{equation}
\end{proposition}
\begin{proof}
With $t=\cos x\in[-1,1]$,
\begin{align*}
\frac{L_1[F](x)}{2\pi e^{-x^2}}
&=1+d^2+3dt+d^2t^2\\
&\ge (1-d)(1-2d)=\frac3{25}>0.
\end{align*}
Indeed, for $d=2/5$ the difference
from $3/25$ is
$(t+1)(4t+26)/25\ge0$.
For $L_2$ take $x=\pi$.
Since $\cos(\pi\pm iy)=-\cosh y$,
\[
F(\pi+iy)F(\pi-iy)
=2\pi e^{-\pi^2}e^{y^2}
(1-d\cosh y)^2.
\]
Its $y^4$ coefficient is
\[
\frac12-\frac{25}{12}d+
\frac{11}{6}d^2
=\frac{22d^2-25d+6}{12}.
\]
At $d=2/5$ this is $-1/25$,
as required.
\end{proof}

\paragraph{Exact tests and formalization scope.}
The standard-library script
\texttt{certificates/finite\_depth\_laguerre\_exact.py}
checks the coefficient formulas by rational
arithmetic through $N=100$, as well as
the exact $d=2/5$ separation.
The Lean module
\texttt{XiKernel/FiniteDepthBarrier.lean}
formally proves the associated
real-algebraic positivity gate and low-order
coefficient identities. A complete Lean
formalization of the Fourier--Gaussian
integral and all-$N$ binomial calculation
is not yet claimed.

\paragraph{Consequence for the Riemann hypothesis.}
This theorem proves, within a positive
analytic kernel class, that even globally
positive inequalities
$L_1,\ldots,L_N$ for any fixed $N$
cannot prove all zeros real.
The point is logical and methodological,
not an assertion of novelty;
Laguerre--P\'olya characterizations are
classical and the correlation viewpoint
has established prior art
\cite{DX}.
For the genuine Riemann $\Xi$, the
research target must therefore be a
uniform argument in \emph{all} orders
$n$, or a materially stronger zero
exclusion principle that explicitly uses
Jacobi/theta arithmetic structure.
